\subsection{Как решать задачи на подсчёт (шпаргалка)}
\textbf{Распознавание.} «Сколько способов» $\to$ спросите: (1) \emph{порядок важен?} (2) \emph{предметы различимы?} (3) \emph{можно ли повторять?} Выбор $k$ из $n$ без порядка -- $\binom nk$; с порядком -- $n^{\underline k}$; с повторами без порядка -- $\binom{n+k-1}k$. Разложение на группы: различимые предметы по различимым ящикам -- $k^n$, по неразличимым -- $S(n,k)$; неразличимые предметы по различимым ящикам -- шарики и перегородки (ниже).
\par \textbf{Приёмы.} (a) \emph{Дополнение:} считать «плохие» и вычесть. (b) \emph{Включения-исключения:} условия «ни одно из $A_i$ не выполнено» $=\sum_{T}(-1)^{|T|}\,|\bigcap_{i\in T}A_i|$; удобно, когда пересечение $T$ условий считается формулой (пример: беспорядки, поверхностные отображения, ограничение $x_i\le k$). (c) \emph{Линейность:} матожидание суммы -- сумма матожиданий, независимость не нужна. (d) \emph{Биекция:} заменить объект на более простой (ПСП $\leftrightarrow$ пути, деревья $\leftrightarrow$ коды Прюфера). (e) \emph{Отражение:} пути, касающиеся запрещённой прямой, биективны путям до отражённой конечной точки.
\subsection{Биномиальные коэффициенты}
\textbf{Базовое.} $\binom nk=\binom n{n-k}=\frac nk\binom{n-1}{k-1}$; Паскаль $\binom nk=\binom{n-1}{k-1}+\binom{n-1}k$; $\binom{-n}{k}=(-1)^k\binom{n+k-1}k$ (отрицательный верхний индекс сводится к положительному).
\par \textbf{Суммы по $k$.} $\sum_k\binom nk=2^n$; чётные и нечётные $k$ дают по $2^{n-1}$; $\sum_k k\binom nk=n2^{n-1}$ (выбрать подмножество и в нём «старшего»); $\sum_k\binom nk\binom km=\binom nm2^{n-m}$ (выбрать $m$-подмножество, потом любое между ним и целым); $\sum_k\binom nk^2=\binom{2n}n$; $\sum_{k}\binom{n-k}k=F_{n+1}$ (разбиения $n$ на единицы и двойки). $\sum_{i\le k}(-1)^i\binom ni=(-1)^k\binom{n-1}k$.
\par \textbf{Вандермонд:} $\sum_k\binom mk\binom n{r-k}=\binom{m+n}r$ -- выбрать $r$ человек из двух групп. \textbf{Клюшка:} $\sum_{i=r}^n\binom ir=\binom{n+1}{r+1}$ (выбор $(r{+}1)$-подмножества по его наибольшему элементу); диагональ $\sum_{i\le k}\binom{n+i}i=\binom{n+k+1}k$.
\par $\sum_{i=0}^k\binom ki\big/\binom ni=\binom{n+1}{n-k+1}\big/\binom nk$. \ $\sum_i\binom ni i^k=\sum_jS(k,j)\,n^{\underline j}\,2^{n-j}$ (степень $i^k$ раскладываем по убывающим степеням).
\par \textbf{Скользящая сумма} $B(l,r,n)=\sum_{i=l}^r\binom ni$: $B(l,r,n{+}1)=2B(l,r,n)+\binom n{l-1}-\binom nr$ -- двигать $n$ на единицу за $O(1)$ (приём для задач «сумма биномов по диапазону при растущем $n$», как в Mo).
\par \textbf{Куммер:} $v_p\binom nm$ равно числу переносов при сложении $m$ и $n-m$ в $p$-ичной системе (показатель простого в биноме). \textbf{Лукас:} $\binom nm\equiv\prod_i\binom{n_i}{m_i}\pmod p$ для простого $p$, где $n_i,m_i$ -- $p$-ичные цифры. Применять, когда $n$ огромно, а $p$ мал.
\par \textbf{Биномы по составному модулю:} раскладываем модуль на степени простых, для каждого считаем $\binom nk$ с учётом $v_p$ (Куммер или Лежандр $v_p(n!)=\sum_i\lfloor n/p^i\rfloor$) и частью без $p$, затем КТО.
\subsection{Разложение на слагаемые (шарики и перегородки)}
Число решений $x_1+\dots+x_n=m$: при $x_i\ge0$ -- $\binom{m+n-1}{n-1}$; при $x_i\ge1$ -- $\binom{m-1}{n-1}$ (вычесть по 1, т.е. «зарезервировать»); при $x_i\ge a_i$ -- вычесть $a_i$ и взять первую формулу; при $0\le x_i\le k$ -- включения-исключения: $\sum_{t=0}^{\min(n,\lfloor m/(k+1)\rfloor)}(-1)^t\binom nt\binom{m-t(k+1)+n-1}{n-1}$ (здесь $t$ -- сколько переменных нарушили границу).
\par $n$ различимых людей по $k$ различимым группам размера $\le m$ (пустые можно): $f(0,k)=1$, $f(n,k)=k\,f(n{-}1,k)-k\binom{n-1}mf(n{-}m{-}1,k{-}1)$.
\par \textbf{Пути на сетке.} Путей $(0,0)\to(a,b)$ шагами вправо/вверх -- $\binom{a+b}a$. Пары путей $(0,0)\to(n,m)$ без общих точек, кроме концов (неупорядоченные): $\binom{n+m-2}{n-1}^2-\binom{n+m-2}n\binom{n+m-2}m$. Обобщение на несколько путей -- Линдстрём--Гессель--Вьенно: det матрицы чисел путей из $A_i$ в $B_j$.
\par \textbf{Задача о голосовании.} За $A$ отдано $a$ голосов, за $B$ -- $b<a$. Порядков, в которых $A$ \emph{строго} впереди в каждый момент: $\frac{a-b}{a+b}\binom{a+b}a$; $A$ \emph{не отстаёт}: $\frac{a-b+1}{a+1}\binom{a+b}a$. Это тот же принцип отражения, что и для Каталана.
\subsection{Каталан и скобочные последовательности}
$C_n=\frac1{n+1}\binom{2n}n=\binom{2n}n-\binom{2n}{n-1}=\frac1{2n+1}\binom{2n+1}n$,\ $C_n=\sum_{i<n}C_iC_{n-1-i}$.\\
{\footnotesize 1, 1, 2, 5, 14, 42, 132, 429, 1430, 4862, 16796, 58786, 208012, 742900, 2674440}.
\par \textbf{Что считает:} правильные скобочные последовательности (ПСП) длины $2n$; триангуляции выпуклого $(n{+}2)$-угольника; бинарные деревья на $n$ вершинах; пути из $(0,0)$ в $(n,n)$, не поднимающиеся выше диагонали; непересекающиеся хорды $2n$ точек на окружности. Если в задаче видите рекурренту $f_n=\sum f_if_{n-1-i}$ -- это Каталан.
\par $\sum_{x_0+\dots+x_k=n}\prod C_{x_i}=\frac{k+1}{n+k+1}\binom{2n+k}n$ (лес из $k{+}1$ таких структур).
\par \textbf{ПСП:} баланс (число «(» минус «)») каждого префикса $\ge0$, итоговый $0$. Скобочных слов с $n$ открывающими и $m\le n$ закрывающими, дополняемых до ПСП: $\binom{n+m}m-\binom{n+m}{m-1}$. Ровно $k$ пиков «()» (числа Нараяны): $N(n,k)=\frac1n\binom nk\binom n{k-1}$; столько же плоских корневых деревьев с $n$ рёбрами и $k$ листьями. Скобка с номером $p$ (с нуля) закрывающая в $\sum_{m=0}^{\min(n-1,\lfloor(p-1)/2\rfloor)}C_mC_{n-m-1}$ ПСП. Идея вывода: из $2n{+}1$ циклических сдвигов слова из $n$ «(» и $n{+}1$ «)» ровно один начинается с ПСП.
\subsection{Числа Стирлинга, Белла, Эйлера}
$S(n,k)$ -- число разбиений $n$-элементного множества на $k$ непустых неупорядоченных блоков. $S(n,k)=kS(n{-}1,k)+S(n{-}1,k{-}1)$ (элемент $n$ либо в одном из $k$ блоков, либо отдельно) $=\frac1{k!}\sum_j(-1)^{k-j}\binom kjj^n$ (включения-исключения по пустым блокам; $k!\,S(n,k)$ -- число сюръекций). Связь со степенью: $x^n=\sum_kS(n,k)\,x^{\underline k}$ -- любая степень раскладывается по убывающим; именно так считают $\sum_i i^n$ и $\sum\binom ni i^n$. Ряд: $S(n,k)=[x^{n-k}]\prod_{i=1}^k\frac1{1-ix}$ (фиксированное $k$); вся строка $n$ одной свёрткой: $S(n,k)=[x^k]\big(\sum_i\frac{i^n}{i!}x^i\big)\big(\sum_j\frac{(-1)^j}{j!}x^j\big)$ (NTT, см. \texttt{opts/poly.cpp}).
\par $c(n,k)$ -- число перестановок $n$ элементов с ровно $k$ циклами (беззнаковые числа I рода): $c(n,k)=c(n{-}1,k{-}1)+(n{-}1)c(n{-}1,k)$. Многочлены: $x^{\overline n}=\sum_kc(n,k)x^k$ и $x^{\underline n}=\sum_k(-1)^{n-k}c(n,k)x^k$. Вся строка $n$ -- коэффициенты $\prod_{t<n}(x+t)$, считаются разделяй-и-властвуй с FFT за $O(n\log^2n)$. Столбец $k$: $\frac{(-\ln(1-x))^k}{k!}$.
\par \textbf{Обращение Стирлинга:} $P(k)=\sum_iS(k,i)Q(i)\iff Q(k)=\sum_i(-1)^{k-i}c(k,i)P(i)$ -- если нашли сумму, выраженную через степени, обратно восстановим коэффициенты по убывающим степеням.
\par Блоки, где любые два элемента одного блока отличаются не менее чем на $d$: $S(n{-}d{+}1,\,k{-}d{+}1)$ при $k\ge d$.
\par \textbf{Белла} $B_n=\sum_kS(n,k)=\sum_k\binom{n-1}{k-1}B_{n-k}$ -- число всех разбиений; {\footnotesize 1, 1, 2, 5, 15, 52, 203, 877, 4140, 21147, 115975, 678570, 4213597}.
\par \textbf{Эйлеровы} $E(n,k)$ -- перестановки с ровно $k$ подъёмами ($p_i<p_{i+1}$): $E(n,k)=(n{-}k)E(n{-}1,k{-}1)+(k{+}1)E(n{-}1,k)=\sum_j(-1)^j\binom{n+1}j(k{+}1{-}j)^n$; $\sum_kE(n,k)=n!$.
\subsection{Перестановки и симметрии}
\textbf{Беспорядки} (без неподвижных точек): $D_n=(n{-}1)(D_{n-1}+D_{n-2})=nD_{n-1}+(-1)^n=n!\sum_{t\le n}\frac{(-1)^t}{t!}=\lfloor n!/e\rceil$ ($n\ge1$); $\sum_k\binom nkD_{n-k}=n!$ (перестановку можно задать числом $k$ неподвижных и беспорядком остальных); {\footnotesize 1, 0, 1, 2, 9, 44, 265, 1854, 14833, 133496}. Типичное применение: «сколько перестановок, где ровно $k$ неподвижных» $=\binom nkD_{n-k}$.
\par \textbf{Средние для случайной перестановки:} циклов -- $H_n$, рекордов -- $H_n$, неподвижных точек -- $1$, инверсий -- $\frac{n(n-1)}4$.
\par Перестановки с циклами только длин из множества $T$: $\sum_n g_T(n)\frac{x^n}{n!}=\exp\big(\sum_{n\in T}\frac{x^n}n\big)$ (EGF).
\par \textbf{Бернсайд:} число классов эквивалентности $=\frac1{|G|}\sum_{g\in G}|X^g|$, где $X^g$ -- объекты, не меняющиеся при $g$. Ожерелья из $n$ бусин, $k$ цветов, с точностью до поворота: $N=\frac1n\sum_{d\mid n}\varphi(d)k^{n/d}$. Браслеты (ещё и отражения): $\frac12\big(N+k^{(n+1)/2}\big)$ при нечётном $n$ и $\frac12\big(N+\frac{(k+1)k^{n/2}}2\big)$ при чётном.
\par Бинарных строк ровно с $m$ различными подпоследовательностями (считая пустую): $\varphi(m{+}1)$.
\subsection{Разбиения числа}
$p(n)$ -- число способов записать $n$ суммой натуральных без учёта порядка. Пентагональная теорема: $\prod_{n\ge1}(1-x^n)=\sum_{k\in\mathbb Z}(-1)^kx^{k(3k-1)/2}$; отсюда $p(n)=\sum_{k\ne0}(-1)^{k+1}p\big(n-\tfrac{k(3k-1)}2\big)$ за $O(n\sqrt n)$ (в сумме только $O(\sqrt n)$ слагаемых). $p(n)\approx\frac1{4\sqrt3n}e^{\pi\sqrt{2n/3}}\approx\frac{0.145}ne^{2.56\sqrt n}$.\\
{\footnotesize $p(0..10)=1,1,2,3,5,7,11,15,22,30,42$;\ $p(20)=627$,\ $p(50)=204226$,\ $p(100)=190569292$}.
\subsection{Деревья, остовы, эйлеровы циклы}
\textbf{Кэли:} помеченных деревьев на $n$ вершинах $n^{n-2}$; код Прюфера -- биекция с последовательностями из $n{-}2$ чисел (повторяя «удали наименьший лист и запиши его соседа»). Деревьев с заданными степенями $d_v$ ($\sum d_v=2n{-}2$, $d_v\ge1$): $\frac{(n-2)!}{\prod(d_v-1)!}$ (вершина $v$ встречается в коде $d_v{-}1$ раз). Лесов из $m$ деревьев на $n$ вершинах, где выделенные вершины $1..m$ лежат в разных деревьях: $m\,n^{n-m-1}$\,$^\dagger$. $\tau(K_{a,b})=a^{b-1}b^{a-1}$.
\par \textbf{Кирхгоф (число остовов):} построить матрицу $L$: для ребра $u\to v$ сделать $L[u][v]{-}{=}1,\ L[v][v]{+}{=}1$ (неориентированное ребро -- оба направления). Минор без строки и столбца $r$ -- число остовных деревьев, направленных \emph{от} $r$ (для неориентированного -- обычных остовов, любой $r$). С весами рёбер -- сумма по остовам произведений весов. Чтобы получить деревья, направленные \emph{к} корню, ставьте $L[u][u]{+}{=}1$ вместо $L[v][v]$.
\par \textbf{BEST:} число эйлеровых циклов орграфа (как циклических последовательностей рёбер) $=t_w(G)\prod_v(\deg_{out}(v){-}1)!$, где $t_w$ -- число остовов, направленных к $w$\,$^\dagger$.
\par \textbf{Эйлеров путь:} в неориентированном -- связность по рёбрам и $\le2$ нечётных вершин; в орграфе -- $\mathrm{out}-\mathrm{in}\in\{0,\pm1\}$, не более одной «начальной» ($+1$) и одной «конечной» ($-1$).
\par \textbf{Простые факты о деревьях:} $\sum_v(\deg v-2)=-2$; листьев не меньше двух; диаметр находится двумя обходами (самая дальняя от любой вершины $a$, затем самая дальняя от неё); центр -- одна или две средние вершины диаметра, радиус $\lceil\mathrm{diam}/2\rceil$.
\subsection{Производящие функции (помеченные объекты, EGF)}
Пара независимых частей -- произведение $A\cdot B$; последовательность -- $\mathrm{Seq}(A)=\frac1{1-A}$; множество -- $\mathrm{Set}(A)=e^{A}$; цикл -- $\mathrm{Cycle}(A)=\ln\frac1{1-A}$. Следствия: $S(n,k)\leftrightarrow\frac{(e^x-1)^k}{k!}$, $c(n,k)\leftrightarrow\frac{(-\ln(1-x))^k}{k!}$, $D_n\leftrightarrow\frac{e^{-x}}{1-x}$, Белл $\leftrightarrow e^{e^x-1}$. Непомеченные объекты: $\mathrm{MSet}(A)=\exp\big(\sum_{k\ge1}\frac{A(z^k)}k\big)$ (разбиения чисел при $A=\frac z{1-z}$), $\mathrm{Cycle}(A)=\sum_{n\ge1}\frac{\varphi(n)}n\ln\frac1{1-A(z^n)}$ (ожерелья).
\subsection{Суммы, ряды, рекуррентности}
$\sum_{i\le n} i^2=\frac{n(n{+}1)(2n{+}1)}6$, \ $\sum i^3=\frac{n^2(n{+}1)^2}4$, \ $\sum i^4=\frac{n(n{+}1)(2n{+}1)(3n^2{+}3n{-}1)}{30}$; \ $\sum_{i=1}^nia^i=\frac{a(na^{n+1}-(n{+}1)a^n+1)}{(a-1)^2}$.
\par $e^x=\sum\frac{x^k}{k!}$, \ $\ln(1{+}x)=\sum_{k\ge1}\frac{(-1)^{k+1}x^k}k$, \ $(1{-}x)^{-n}=\sum_k\binom{n+k-1}{n-1}x^k$, \ $\arctan x=\sum(-1)^k\frac{x^{2k+1}}{2k+1}$.
\par \textbf{Линейная рекуррента} $a_n=\sum_{i\le k}c_ia_{n-i}$: решение -- $a_n=\sum d_jr_j^n$ по корням $r_j$ характеристического многочлена $x^k-\sum c_ix^{k-i}$ (кратный корень даёт множитель-многочлен). Практически: $a_n$ за $O(k^2\log n)$ возведением $x$ в степень по модулю характеристического многочлена; по $2k$ первым членам рекурренту восстанавливает Берлекэмп--Мэсси (\texttt{math/linrec.cpp}).
